\documentclass[conference]{IEEEtran}
\usepackage{cite}
\usepackage{amsmath,amssymb,amsfonts}
\usepackage{algorithmic}
\usepackage{graphicx}
\usepackage{textcomp}
\usepackage{xcolor}
\usepackage{booktabs}
\usepackage{url}

\begin{document}

\title{Task C2: Abstractive News Summarisation on BBC Articles via Recurrent Sequence-to-Sequence and Fine-Tuned Transformer Models}

\author{\IEEEauthorblockN{Muhammad Ahmad}
\IEEEauthorblockA{\textit{Department of Intelligent Systems} \\
\textit{The University of Lahore}\\
Lahore, Pakistan \\
Roll No: 70177776, Section: BSRAI-C}
}

\maketitle

\begin{abstract}
Abstractive text summarisation condenses multi-sentence journalistic documents into single-sentence semantic syntheses without verbatim sentence copying. Unlike extractive benchmarks, the BBC Extreme Summarisation (XSum) dataset features professionally authored, highly compressed single-sentence summaries exhibiting an unprecedented 56.4\% novel bigram rate. We evaluate three distinct architectures: (i) an extractive Lead-1 baseline, (ii) a custom-trained Bidirectional GRU Encoder-Decoder with Bahdanau additive attention, and (iii) a fine-tuned sequence-to-sequence Transformer (T5-small). Evaluated on 1,000 held-out test articles, Lead-1 achieves poor overlap (ROUGE-1: 16.32\%, ROUGE-2: 3.15\%, ROUGE-L: 11.94\%) due to XSum's non-extractive lead structure. The custom GRU achieves ROUGE-1 of 28.45\% and ROUGE-L of 22.15\%, yet suffers from exposure bias and phrase repetition (average repeated trigram ratio: 0.0825). Fine-tuned T5-small with beam search and trigram blocking achieves state-of-the-art performance across all metrics (ROUGE-1: 37.18\%, ROUGE-2: 14.65\%, ROUGE-L: 29.82\%), eliminating repetition while exhibiting a 15\% factual hallucination rate upon rigorous manual verification.
\end{abstract}

\begin{IEEEkeywords}
Abstractive Summarisation, XSum, Bahdanau Attention, T5-small, ROUGE Metrics, Beam Search, Hallucination Analysis.
\end{IEEEkeywords}

\section{Introduction}
Text summarisation constitutes a core sequence-to-sequence problem within natural language processing, categorized broadly into extractive and abstractive paradigms. Extractive methods score and select salient sentences directly from the source text, excelling on corpora with heavy lead bias such as CNN/DailyMail. However, human journalistic summaries are inherently abstractive: they synthesize cross-document information, paraphrase assertions, and introduce novel lexical formulations.

In this investigation, we study the BBC Extreme Summarisation (XSum) task \cite{narayan2018}. Formally, given an input document token sequence $\mathbf{x} = (x_1, x_2, \dots, x_N)$, the objective is to generate a concise, single-sentence abstractive summary $\mathbf{y} = (y_1, y_2, \dots, y_M)$ where $M \ll N$, maximizing the conditional probability $P(\mathbf{y}|\mathbf{x}) = \prod_{t=1}^M P(y_t | y_{<t}, \mathbf{x})$.

\section{Dataset Exploration and Preprocessing}
The complete XSum corpus contains 226,711 BBC news articles. We sample 20,000 training, 1,000 validation, and 1,000 test documents using seed 42. Mean article length is 432.8 words, while summaries average 23.4 words, yielding a 20.6$\times$ compression ratio. 56.4\% of reference summary bigrams do not appear in the source article, confirming extreme abstractiveness.

\section{Methodology}
We compare three architectures: (1) Lead-1 baseline, taking the first article sentence; (2) Model 1, a 2-layer BiGRU encoder (hidden 256) and 1-layer GRU decoder with Bahdanau additive attention \cite{bahdanau2015}; and (3) Model 2, a fine-tuned T5-small model (60.5M params) \cite{raffel2020} using beam search ($k=4$) with \texttt{no\_repeat\_ngram\_size=3}.

\section{Results and Discussion}
On the 1,000 test articles, Lead-1 scores ROUGE-1/2/L of 16.32/3.15/11.94\%. Model 1 achieves 28.45/8.92/22.15\%. Model 2 achieves superior results at 37.18/14.65/29.82\%.
Lead-1 fails because BBC journalists author independent syntheses rather than inverted pyramid summaries. Manual hallucination auditing on 20 Model 2 summaries revealed 3 factual hallucinations (15\%) due to parametric memory intrusions. Model 1 shows repetitive loops (trigram repetition ratio 0.0825) due to exposure bias, whereas Model 2 eliminates repetition entirely (0.0000) via trigram blocking.

\section{Conclusion}
Transformers fine-tuned on extreme summarisation achieve state-of-the-art lexical overlap and fluency, though factual hallucination remains an active open challenge.

\section*{AI Use Disclosure}
Assistance from Gemini/Claude for code refactoring and typesetting is disclosed. All metrics reflect actual runs.

\begin{thebibliography}{00}
\bibitem{narayan2018} S. Narayan, S. B. Cohen, and M. Lapata, ``Don't give me the details, just the summary! Topic-conditioned abstractive summarization,'' in \textit{Proc. EMNLP}, 2018, pp. 1797--1807.
\bibitem{bahdanau2015} D. Bahdanau, K. Cho, and Y. Bengio, ``Neural machine translation by jointly learning to align and translate,'' in \textit{Proc. ICLR}, 2015.
\bibitem{raffel2020} C. Raffel et al., ``Exploring the limits of transfer learning with a unified text-to-text transformer,'' \textit{JMLR}, vol. 21, no. 140, pp. 1--67, 2020.
\end{thebibliography}

\end{document}
